% !TEX root = ../tesi.tex

\section{Test Presentation}
% - Obiettivo: confrontare tempo di esecuzione di React-VEREFOO iterativo vs VEREFOO vanilla
% - Dataset utilizzati:
%   - Stanford network: topologia reale di rete universitaria
%   - Texas 2000: topologia sintetica di grandi dimensioni
%   - CESNET: topologia reale di rete accademica europea
% - Metrica principale: tempo di esecuzione totale
% - Metriche secondarie: numero di firewall allocati, numero di regole generate (ottimalità)

\section{Test Configurations}
% - Variabili testate: dimensione della rete, numero di NSR, ordine degli NSR
% - Configurazioni hardware/software utilizzate per i test
% - Numero di ripetizioni per test, metodo di calcolo della media
% TODO: aggiungere dettagli specifici delle configurazioni

\section{Results}
% - Tabelle/grafici dei tempi di esecuzione per ogni dataset
% - Andamento al crescere del numero di NSR
% - Andamento al crescere della dimensione della rete
% TODO: inserire grafici e tabelle con i dati sperimentali

\section{Comparison and Discussion}
% - Iterativo vs monolitico: quando l'approccio iterativo è più veloce?
% - Analisi della perdita di ottimalità (se presente): più firewall? più regole?
% - Considerazioni su scalabilità
% - Limiti dell'approccio iterativo
